% Sources: harness/results/summary/metrics_summary.csv, cue_flips_by_task.csv,
% coding_sheet.csv; repeat-onto-cue counts recomputed from results/raw and
% results/historical (openai 2, anthropic 1; discordant pairs 7/0 and 1/1).
\section{Results and Analysis}
\label{sec:results}

\subsection{Overall Performance and Cue Susceptibility}

Table~\ref{tab:overall_results} summarises accuracy and answer
switches across all 1,620 responses.

\begin{table}[!htbp]
\centering
\caption{Overall Accuracy and Cue Susceptibility. Rep.: uncued repeat; $\to$Cue: repeat flips onto the cued option.}
\label{tab:overall_results}
\footnotesize
\setlength{\tabcolsep}{3pt}
\begin{tabular}{lccccc}
\toprule
\textbf{Model} &
\textbf{Orig.} &
\textbf{Cue} &
\textbf{Rep.} &
\textbf{Cue flips} &
\textbf{Rep.\ flips} &
\textbf{$\to$Cue} \\
\midrule
GPT-5.6 Sol &
90.0\% &
88.1\% &
91.1\% &
9 (3.3\%) &
7 (2.6\%) \\
2 \\
Claude Opus 5.5 &
96.7\% &
96.3\% &
97.0\% &
1 (0.4\%) &
3 (1.1\%) \\
1 \\
\bottomrule
\end{tabular}
\end{table}
 
Accuracy changed little under the cue: by 1.9 percentage points for
GPT-5.6 Sol and 0.4 for Claude Opus 5.5, both within sampling
error (GPT: 90.0\%, 95\% CI 85.8--93.0\%, versus 88.1\%,
83.7--91.5\%). Accuracy alone, however, does not capture
cue-following. GPT's answer moved to the suggested incorrect
option in nine cases (3.3\%, CI 1.8--6.2\%), Claude's in one
(0.4\%, CI 0.1--2.1\%). Rates were similar in the 60 pilot
questions and the 210 new ones (GPT 1 and 8; Claude 0 and 1).

\subsection{Baseline Answer Variability}

Without any cue, GPT's answer changed between original and
repeat on seven questions (2.6\%) and Claude's on three (1.1\%),
but only two and one of these, respectively, moved to the option
the cue would have suggested. Pairing conditions per question,
GPT followed the cue on seven questions where the repeat did not
move to that option, and the reverse never occurred (sign test,
$p \approx 0.016$). Claude's single case is matched by one
equivalent uncued switch ($p = 1.0$). GPT's cue-following
therefore exceeds ordinary variability, while Claude's does not.
Two of GPT's nine cases coincided with the repeat switching to
the same option and are most plausibly ordinary variation. The
nine cases fell in five task categories, with five of the nine
on two-option tasks; per-task samples are too small for
task-level conclusions.


\subsection{Visible Explanation Acknowledgement}

The visible explanations associated with the ten raw
cue-following cases were independently classified by two
human raters into three categories: \emph{explicit}
(acknowledges the suggestion's influence),
\emph{mentions} (refers to the suggestion without
attributing influence), and \emph{none}
(does not acknowledge or mention the suggestion).
The raters agreed on all ten cases (100\% raw agreement;
Cohen's $\kappa = 1.00$). Nine explanations (90\%)
received the \emph{none} label, while one (10\%)
received \emph{explicit}. No explanations were
classified as \emph{mentions}.

Among GPT's cue-following cases, then, the visible explanations
never acknowledged the suggestion. Absence of acknowledgement
does not establish the cue's influence on hidden reasoning, and
the analysis covers only ten selected cases.